\section{Twistor spaces of four- and six-manifolds}\label{section3}

In this section we describe the integral cohomology of the twistor space of (Riemannian) almost complex six-manifolds, along with the integral cohomology of the twistor space of arbitrary oriented (Riemannian) four-manifolds. Finally we compute the Chern classes of the Atiyah--Hitchin--Singer almost complex structure on the latter.

\subsection{Spinors on almost complex four- and six-manifolds}

Let $M$ be an oriented Riemannian manifold and let $Z(M)$ denote the space of orthogonal almost complex structures on $M$, called the \emph{twistor space}. The space
$Z(M)$ has two components corresponding to the structures which induce the given orientation on $M$ or the opposite orientation. Recall, we
denote these by $Z_{\pm}(M)$, and refer to them as the \emph{positive} and \emph{negative} twistor spaces.

The Atiyah--Hitchin--Singer almost complex structure on $Z(M)$ is defined as follows: the Levi-Civita connection induces a connection on the bundle $Z(M) \to M$ splitting $TZ(M)$ into a~direct sum of horizontal and vertical subspaces
\[
TZ(M) =T^{\rm h} Z(M) \oplus T^{\rm v}Z(M).
\]
The almost complex structure on $T^{\rm v} Z(M)$ is determined by the algebraic variety structure of the fibers $J(T_x M)$ of $Z(M) \to M$ explained in Section~\ref{section2}. At a point $J_x$ in the fiber
over $x\in M$ the complex structure on $T^{\rm h}_{J_x} Z(M) = T_x M$ is given\footnote{The Eells--Salamon almost complex structure is obtained replacing the vertical component of the Atiyah--Hitchin--Singer almost complex structure by its negative; see~\cite[Section~3]{Sa84} for details.} by $J_x$.



A ${\rm spin}^{\rm c}$ structure on $M$~\cite[Appendix~D]{LM} yields a complex spinor
bundle $S_\C(M)\to M$ together with an action
\[
\CCl(TM) \otimes S_\C(M) \to S_\C(M).
\]
The discussion in the previous section then provides a one-to-one correspondence between projectivized pure spinors on $M$ and orthogonal complex structures inducing a given orientation:
\[
Z_{\pm}(M) \cong \mathbb{P}\big( {\mathcal P}S _{\C}^{\pm}(M)\big).
\]
Now suppose $M$ is given an orthogonal almost complex structure $J$ compatible with the orientation. Then $J$ gives rise to a canonical ${\rm spin}^{\rm c}$ structure on $M$ via the group homomomorphism ${\rm U}(n) \xrightarrow{\kappa} \Spin^{\rm c}(2n)$, which is the unique
lift as a group homomorphism in the following diagram{\samepage
\[
\begin{tikzcd}
& \Spin^{\rm c}(2n) = \big(\Spin(2n) \times S^1\big)/\{\pm(1,1)\} \ar[d, "\pi \times \delta"] \\
{\rm U}(n) \ar[ur,"\kappa"] \ar[r,"\iota \times \det"] & {\rm SO}(2n) \times S^1.
\end{tikzcd}
\]
Here $\iota$ denotes the inclusion, $\pi\big(\big[h, {\rm e}^{{\rm i}\theta}\big]\big) = [h]$, and $\delta\big(\big[h,{\rm e}^{{\rm i}\theta}\big]\big)={\rm e}^{2{\rm i}\theta}$.}

An element $g\in {\rm U}(n)$ is diagonalized by some orthonormal basis $(e_1,Je_1, \ldots, e_n, Je_n)$. If
\[g=\mathrm{diag}\big({\rm e}^{{\rm i}\theta_1},\ldots, {\rm e}^{{\rm i}\theta_n}\big),\]
then
\[
\kappa(g)= \left[ \prod_{k=1}^n \left( \cos \frac{\theta_k} 2 + {\rm i} \sin \frac{\theta_k}{2} e_k J e_k\right), {\rm e}^{{\rm i} \frac{\sum \theta_k}{2}} \right]
\]
 (cf.~\cite[equation~(D.10)]{LM}). Using
the formulas for the action of $e_i Je_i$ in the proof of Lemma~\ref{plusmin}, we see that the spinor action of $\kappa(g)$
on $\Lambda^*_{\C}\big(\R^{2n}\big)$, where $\R^{2n}$ is equipped with the complex structure given by $\mathrm{diag}\left( \left(\begin{smallmatrix} 0 & -1 \\ 1 & 0 \end{smallmatrix}\right), \ldots, \left(\begin{smallmatrix} 0 & -1 \\ 1 & 0 \end{smallmatrix}\right) \right)$, agrees with the standard action of $g$ on this space. We obtain the following result:

\begin{Proposition}\label{identification}
Let $M$ be an oriented Riemannian manifold with an orthogonal almost complex structure $J$ compatible with the orientation. Then we have the following isomorphisms of smooth fiber bundles over $M$:
\begin{enumerate}\itemsep=0pt
\item[$1.$] If $\dim(M)=4$, then
\[
Z_+(M) \cong \mathbb{P}\big( \C \oplus \Lambda^2_\C(TM)\big), \qquad Z_{-}(M) \cong \mathbb{P}( TM).
\]
\item[$2.$] If $\dim(M)=6$, then
\[
Z_+(M) \cong \mathbb{P}\big( \C \oplus \Lambda^2_\C(TM)\big), \qquad Z_{-}(M) \cong \mathbb{P}\big( TM \oplus \Lambda_\C^3(TM)\big).
\]
\end{enumerate}
\end{Proposition}

\begin{Remark}
 The Atiyah--Hitchin--Singer and Eells--Salamon almost complex structures on (either) twistor space do not agree with the natural almost complex structure on the projectivization. Indeed, the projection map from the projectivization is complex linear whilst this is not the case for the twistor projection.
\end{Remark}

The projective bundle formula now gives the following formulas for the cohomology rings of the components of the twistor space in terms of the
Chern classes of $(TM, J)$.

\begin{Corollary}
\label{cohomtwis}
Let $M$ be an oriented Riemannian manifold with an orthogonal almost complex structure $J$ compatible with the orientation. We have the following isomorphisms as algebras over $H^*(M)$:
\begin{enumerate}\itemsep=0pt
\item[$1.$] If $\dim(M)=4$, then
\begin{align*}
&H^*(Z_+(M))\cong H^*(M)[x]/\big(x^2+c_1 x\big), \\
&H^*(Z_{-}(M))\cong H^*(M)[x]/\big(x^2 + c_1 x + c_2\big).
\end{align*}
\item[$2.$] If $\dim(M)=6$, then
\begin{align*}
&H^*(Z_+(M))\cong H^*(M)[x]/\big(x^4 +2c_1x^3 +\big(c_1^2+c_2\big) x^2 + (c_1c_2-c_3)x\big), \\
&H^*(Z_{-}(M))\cong H^*(M)[x]/\big(x^4 + 2c_1x^3+ \big(c_1^2+c_2\big)x^2 + (c_1c_2+c_3)x\big). \end{align*}
\end{enumerate}
\end{Corollary}
